
For the first step, one must transform @@M0@@ into the identity matrix @@M1@@ in a Zariski neighborhood of @@M2@@ by multiplying @@M3@@ on the left and on the right by suitable matrices @@M4@@. Let us first show that one can thereby make @@M5@@. If @@M6@@ and @@M7@@, we seek a matrix @@M8@@ such that the product @@M9@@ has @@M10@@. Since @@M11@@, there is at least one index @@M12@@ such that @@M13@@. Then one has @@M14@@, and @@M15@@ answers the question. One may therefore suppose @@M16@@ and, after localizing, suppose that @@M17@@ is invertible. If @@M18@@, then in @@M19@@ one has
@@M20@@. After localizing, one may then suppose that @@M21@@ is invertible. Then @@M22@@ is invertible and one finds @@M23@@ in @@M24@@.

Supposing @@M25@@, one can then put @@M26@@ in the form
@@M27@@
in a Zariski neighborhood, of @@M28@@ by multiplying @@M29@@ on the left (resp. on the right) by matrices @@M30@@ (resp. @@M31@@). Next one makes @@M32@@, etc. By induction, one can then transform @@M33@@ in a Zariski neighborhood @@M34@@ of @@M35@@ into a matrix @@M36@@ with @@M37@@ for @@M38@@ or @@M39@@. But by the very construction one then has @@M40@@ on @@M41@@, where @@M42@@ and @@M43@@ are products of matrices @@M44@@. Thus one has
@@M45@@
so that @@M46@@, @@M47@@, and the result follows.

We therefore suppose that @@M48@@ is already, on @@M49@@, a product of matrices @@M50@@.
It remains to see that these are products of commutators in an étale neighborhood of @@M51@@. First, one sees that any two matrices
